\chapter{Diseño del sistema}

\section{Diseño general}
Tech-Flow es una aplicación de escritorio de una sola ventana (\emph{ventana única}) con tres roles. Cada rol tiene su propia vista FXML; el controlador de esa vista orquesta la pantalla, las clases de dominio ejecutan las reglas y \texttt{MotorCSV} persiste en archivos. No hay servidor, ni red, ni base de datos: la unidad de consistencia es el archivo.

\section{Arquitectura y organización de módulos}
\begin{table}[H]\centering\caption{Módulos principales y responsabilidades}\begin{tabular}{p{.26\textwidth}p{.6\textwidth}}\toprule
Módulo & Responsabilidad\\\midrule
\texttt{com.overbyte} & Aplicación JavaFX: \texttt{Navegador} (ventana única, transiciones, overlays) y \texttt{Main} (entrada).\\
\texttt{com.overbyte.controllers} & Controladores por pantalla; ayudantes \texttt{AyudasUI}, \texttt{InformeTexto}, \texttt{AyudasArchivo}, \texttt{FiltrosCategoria}.\\
\texttt{user} & Dominio: \texttt{Empleado} (base), \texttt{Gestor}, \texttt{Vendedor}, \texttt{Almacen}, \texttt{Rol}.\\
\texttt{models} & Entidades: \texttt{Producto}, \texttt{Proveedor}, \texttt{Cliente}, \texttt{Inventario}.\\
\texttt{services} & \texttt{MotorCSV} (parseo/escritura), \texttt{Verificador} (autenticación), \texttt{ValidadorDatos}.\\
\texttt{config} & \texttt{ArchivosCSV} (nombres canónicos), \texttt{ConfiguracionCSV}.\\
\texttt{com.overbyte.vista} & FXML: \texttt{login}, \texttt{admin}, \texttt{vendedor}, \texttt{almacen}.\\
\texttt{com.overbyte.styles} & \texttt{estilos.css} con las clases del sistema.\\\bottomrule\end{tabular}\end{table}

\section{Diseño de interfaces}
La ventana única se construye en \texttt{Navegador.iniciar(stage)}: un \texttt{StackPane} raíz con dos capas, la de la vista activa y la de overlays (avisos, formularios y confirmaciones). El cambio de pantalla es con fundido y deslizamiento corto. Las interacciones entre vistas y dominio: la vista recibe un \texttt{Empleado} autenticado mediante \texttt{setEmpleado(...)} y con él instancia el servicio correspondiente (\texttt{Vendedor}, \texttt{Gestor} o \texttt{Almacen}).

\section{Diseño de los procesos principales}
\subsection{Venta confirmada}
\begin{enumerate}
\item El cajero busca productos (filtro opcional por categoría) y arma el carrito; el carrito valida stock por línea.
\item Al confirmar, se exige nombre, CI y teléfono del comprador (\texttt{ValidadorDatos}).
\item \texttt{generarOrdenDeVenta(...)} escribe líneas \emph{PROFORMA} en \texttt{ventas.csv} y releyendo el archivo antes de persistir.
\item \texttt{confirmarVenta(...)} descuenta existencias en \texttt{inventario.csv}, registra/actualiza el cliente en \texttt{clientes.csv}, crea garantías por línea vendida en \texttt{garantias.csv} y marca las filas \emph{CONFIRMADA}.
\item La GUI presenta el resumen en un cuadro interno con opción de guardar el comprobante .txt.
\end{enumerate}

\subsection{Apertura de caja}
El administrador confirma la apertura en un cuadro interno; el estado queda en la única fila de \texttt{caja.csv}. Sin estado \emph{ABIERTA}, \texttt{confirmarVenta} rechaza la operación: la caja es el interruptor operativo del punto de venta.

\section{Diseño de persistencia y entidades}
\begin{table}[H]\centering\caption{Esquema de los archivos CSV}\begin{tabular}{p{.24\textwidth}p{.64\textwidth}}\toprule
Archivo y columnas\\\midrule
\texttt{empleado.csv} & ci, nombre, rol, contrasena\\
\texttt{cliente.csv/clientes.csv} & idCliente, nombreCliente, ci, telefono\\
\texttt{inventario.csv} & idProducto, nombre, marca, categoria, descripcion, precioEntrada, precioVenta, stock\\
\texttt{ventas.csv} & idVenta, idOrden, estado, ciCliente, nombreCliente, idProducto, nombreProducto, cantidad, precioUnitario, total, fechaVenta, ciEmpleado, idCliente, telefonoCliente, idGarantia\\
\texttt{garantias.csv} & idGarantia, idVenta, idOrden, idCliente, ciCliente, nombreCliente, telefonoCliente, idProducto, nombreProducto, fechaInicio, fechaFin, estado, ciEmpleado\\
\texttt{caja.csv} & estado, fechaApertura, fechaCierre, ciEmpleado\\
\texttt{entradas.csv} & idEntrada, idProveedor, idProducto, nombreProducto, cantidad, precio, fechaEntrada, ciEmpleado\\
\texttt{proveedores.csv} & idProveedor, nombreProveedor, contactoProveedor\\
\texttt{pedidos.csv} & idPedido, idProveedor, idProducto, cantidad, fechaPedido, ciEmpleado\\\bottomrule\end{tabular}\end{table}

Relaciones: una venta genera una garantía por idProducto vendido; idCliente en \texttt{ventas} y \texttt{garantias} referencia \texttt{clientes.csv}; \texttt{entradas} referencia \texttt{productos} y \texttt{proveedores} cuando aplica.

\section{Estructuras de datos y complejidad}
\begin{itemize}
\item \textbf{Carrito} (\texttt{LinkedHashMap<idProducto, cantidad>}): inserción, totales y vaciado en $O(1)$ amortizado; orden de inserción conservado en el recibo.
\item \textbf{Búsqueda de inventario}: \texttt{buscarEnInventario} recorre el CSV — $O(n)$ por consulta; con $n\approx 5\cdot10^2$ productos el coste es irrelevante.
\item \textbf{Tablas de ventas/entradas/inventario}: reconstrucción completa por render en $O(n)$; sin paginación ni virtualización.
\item \textbf{Autenticación}: comparación en texto plano sobre $e$ filas; sin dependencias externas.
\end{itemize}
